\notechapter{N-20221228}{Quotient Topology and Product Topology}

\section{Interval modulo endpoints}

The note begins with the quotient
\[
  D^1=[-1,1],\qquad -1\sim1.
\]
Identifying the two endpoints gives a circle:
\[
  D^1/(-1\sim1)\cong S^1.
\]
A small neighborhood around the glued point corresponds to two half-interval neighborhoods at \(-1\) and \(1\) upstairs.  This is exactly the local reason quotient topology is defined by preimages.

More generally,
\[
  D^n/S^{n-1}\cong S^n.
\]
Collapsing the boundary of a disk to one point produces a sphere.

\section{More quotient spaces}

The note records the standard square identifications.  If a square has opposite vertical sides identified in the same direction, one obtains a cylinder.  If the two circular boundary components of that cylinder are then identified in the same direction, one obtains a torus.  If one pair of edges is identified with a twist, one obtains a Möbius band or a Klein bottle depending on the full gluing pattern.

Another example is
\[
  \mathbb R/\mathbb Z\cong S^1,
\]
where \(x\sim y\) iff \(y-x\in\mathbb Z\).  This is the real line wrapped around the circle.

\section{Definition of quotient topology}

Let \(q:X\to X/{\sim}\) be the quotient map.  A set \(U\subseteq X/{\sim}\) is open iff
\[
  q^{-1}(U)\text{ is open in }X.
\]
This is the finest topology on the quotient making \(q\) continuous.  The note emphasizes that quotient topology is not guessed by drawing; it is forced by the universal property of the quotient map.

\section{Product topology}

For spaces \(X,Y\), the product topology on \(X\times Y\) is generated by basic open rectangles
\[
  U\times V,\qquad U\subseteq X\text{ open},\quad V\subseteq Y\text{ open}.
\]
Thus a set is open in \(X\times Y\) if it is a union of such rectangles.

The page's picture of a small rectangle in a grid illustrates the basic rule: local neighborhoods in a product are products of local neighborhoods in the factors.


\section{Universal property of quotient topology}

The quotient topology is not chosen by visual intuition alone.  Let \(q:X\to Y=X/{\sim}\) be the quotient map.  A map \(f:Y\to Z\) is continuous iff
\[
  f\circ q:X\to Z
\]
is continuous.  This is the universal property of the quotient topology.

This property explains the definition.  We declare \(U\subseteq Y\) open precisely when \(q^{-1}(U)\) is open in \(X\), because this is exactly what makes continuity out of the quotient detectable before quotienting.

\section{Endpoint identification in detail}

For \([0,1]/(0\sim1)\), points near the glued point come from two pieces upstairs: a small interval near \(0\) and a small interval near \(1\).  This is why the quotient looks like a circle rather than an interval with a strange endpoint.  The local neighborhoods have been glued.

Similarly, \(D^2/S^1\cong S^2\): collapse the boundary circle of a disk to a single point.  The interior of the disk becomes the sphere with one point removed, and the collapsed boundary becomes the missing point.

\section{Product topology via projections}

The product topology on \(X\times Y\) is the coarsest topology making the projection maps
\[
  \pi_X:X\times Y\to X,
  \qquad
  \pi_Y:X\times Y\to Y
\]
continuous.  Its basis consists of rectangles
\[
  U\times V
\]
where \(U\subseteq X\) and \(V\subseteq Y\) are open.

The word “coarsest” matters.  We do not add more open sets than necessary.  Once all projection preimages \(\pi_X^{-1}(U)=U\times Y\) and \(\pi_Y^{-1}(V)=X\times V\) are open, finite intersections give \(U\times V\), and arbitrary unions give the full product topology.

\section{Quotient and product as opposite directions}

Product topology builds a larger space from two spaces while preserving maps out to the factors.  Quotient topology builds a smaller space by identifying points while preserving maps out of the quotient.  Product is controlled by projections; quotient is controlled by the quotient map.

This contrast is a useful memory:
\[
  \text{product: maps out to factors},\qquad
  \text{quotient: maps out after gluing}.
\]
Both definitions are forced by continuity requirements rather than by pictures.

\section{A second lens}

Quotient topology and product topology are dual-looking constructions.  A quotient identifies points and defines open sets by pulling back along the quotient map.  A product combines spaces and defines open sets by rectangles from the two factors.  Both are examples of topology being controlled by maps rather than by drawings alone.


\paragraph{Expanded synthesis.}
For this chapter, the elementary material should be read as a study of what remains invariant when coordinates change. The earlier sections (`Interval modulo endpoints`, `More quotient spaces`, `Definition of quotient topology`, `Product topology`, `Universal property of quotient topology`) compute with arrays, bases, pivots, determinants, or eigenvectors; the deeper object is a linear map together with the spaces it acts on. A matrix is a representation of that map after bases have been chosen. This is why formulas such as
\[
  A' = P^{-1}AP,\qquad \widetilde A = T^{-1}AS
\]
are not cosmetic: they distinguish an intrinsic morphism from a coordinate description.

The structural hierarchy is as follows. Kernels record what a map forgets, images record what it reaches, cokernels record obstructions to solvability, and quotients record information after an equivalence has been imposed. Determinants act on the top exterior power and therefore measure oriented volume; traces are invariant under cyclic rearrangement and become characters in representation theory; adjoints depend on an inner product and lead to self-adjoint spectral theory.

A useful way to return to the computations is to ask, for every formula, which invariant it is measuring. Row reduction measures rank and kernel; diagonalization measures decomposition into invariant subspaces; Gram--Schmidt measures orthogonal projection; positive definiteness measures ellipsoid geometry; tensor notation measures how components transform. The elementary methods are therefore the finite-dimensional laboratory in which duality, quotienting, functoriality, and spectral decomposition can be seen clearly.
